\clearpage
\section{Especificação formal M11: turma, matrícula e convite}
\label{sec:formal-m11}
\sloppy

M11 formaliza o alvo normativo de turma, matrícula ativa, vínculo canônico
Aluno--Turma, espelho de consulta e convite de ingresso (RF17, RF18, RN-TUR-01,
UI-10), incluindo a decisão humana HQ-M11-001 = A. A evidência cobre somente o
modelo formal declarado e não homologa a implementação Firebase.

\subsection*{CUE}
O contrato CUE (\texttt{\#M11Contrato}) verifica shapes, enums fechados,
obrigatoriedade, nulabilidade e limites locais de turma, vínculo canônico,
espelho, evento de inclusão/exclusão, convite, chave de pendência e resultado de
aceitação. Condicionais são verificadas: justificativa obrigatória exatamente
quando a exceção é declarada, autoria/instante do aceite no estado aceito, modo
de ingresso na inclusão e contexto \texttt{GLOBAL} quando \texttt{id\_turma} é
nulo. O CUE não prova concorrência temporal: unicidade, transições, contador,
exceção e composição pertencem ao Alloy; a forma canônica do e-mail e as regras
determinísticas de capacidade pertencem ao Rust.

\subsection*{Identidade e ciclo de vida da turma}
\texttt{Turma} tem dono (\texttt{id\_professor}), matéria existente, nome, ano,
semestre 1/2, capacidade positiva, \texttt{codigo\_turma} único gerado no
servidor e \texttt{status} Ativo/Arquivada. \texttt{codigo\_turma} é reservado
permanentemente em \texttt{Chaves\_Unicas/Turma\_\_\{codigo\}} e arquivar não o
libera para reuso. Arquivar/desarquivar incrementam \texttt{versao} e preservam
ID, membros, código e histórico; a turma arquivada bloqueia ingresso e aceite de
convite de turma (somente leitura), e o desarquivamento restaura Ativo com a
mesma identidade.

\subsection*{Matrícula canônica, contador e espelho}
O vínculo \texttt{Turma/\{id\}/Alunos/\{uid\}} é a autoridade de matrícula; o
espelho \texttt{Usuarios/\{uid\}/Turmas/\{id\}} é modelado como projeção exata
do vínculo canônico, portanto não sobrevive à remoção e nunca autoriza sozinho.
O contador \texttt{qtd\_alunos} é igual à cardinalidade dos vínculos canônicos
ao fim de toda transição e nunca é negativo. Inclusão e remoção produzem
exatamente um evento histórico cada, e a remoção preserva os eventos anteriores.
Removido anteriormente não reingressa por código; o reingresso exige convite.

\subsection*{Capacidade, exceção e HQ-M11-001 = A}
O ingresso ordinário exige \texttt{qtd\_alunos < capacidade} no commit; sem
vaga, nenhuma transição ordinária existe. A exceção nominal válida — convite
pendente do professor dono com \texttt{exceder\_capacidade} e justificativa
persistida — pode ultrapassar a capacidade e deixar \texttt{qtd\_alunos >
capacidade} sem invalidar a turma; não é bypass genérico e o código de turma
nunca autoriza exceção. Editar/reduzir a capacidade para valor abaixo da
ocupação preexistente é proibido, fail-closed (HQ-M11-001 = A).

\subsection*{Convites, pendência e idempotência}
A chave determinística de pendência (HMAC de e-mail normalizado e turma ou
GLOBAL) é distinta do ID imutável do convite. Há no máximo uma pendência por
(e-mail, contexto); aceitação ou expiração a liberam. Um novo convite após
terminalidade preserva o convite aceito/expirado e sua auditoria; o reenvio de
convite pendente mantém o mesmo documento. Aceitar convite de turma cria
vínculo, espelho, evento e incrementa o contador; aceitar convite global não
cria matrícula. A aceitação consome o convite uma vez; retry do mesmo ator,
token e operação devolve o receipt (M7), e reuso incompatível não o herda.

\subsection*{Composição M7 e M9}
A composição com M7 é abstrata: retry não duplica vínculo, evento nem contador,
e um mesmo \texttt{cmdId} exige a mesma identidade de comando. A composição com
M9 assume \texttt{authOk} (usuário ativo e versão de permissões corrente):
gerir turma exige o professor dono (ou a intervenção Q13 da chefia), a revogação
impede commit posterior e a leitura exige vínculo canônico atual, dono ou
chefia. Autorização é necessária, nunca suficiente.

\subsection*{Responsabilidade das camadas}
CUE valida shapes e limites locais. Alloy valida unicidade de código, ciclo de
vida, ingresso por vaga, exceção acima da capacidade, proibição de reduzir
capacidade abaixo da ocupação, coerência vínculo--contador--espelho, remoção e
reingresso, convites e composição M7/M9. Rust valida proveniência, hashes,
conjunto ordenado do receipt e as regras determinísticas de ingresso, exceção e
canonicalização de e-mail. O guard mecânico compara enums, nomes e predicados
centrais entre documentação, CUE, Alloy e Rust.

\begingroup\small\raggedright
% Gerado por lcqui-spec-doc; não editar.
\subsection*{M11\_Turma\_Matricula\_Convite --- formalização executável M11}
Shapes estruturais de turma, vínculo canônico Aluno--Turma, espelho mínimo de consulta, evento de inclusão/exclusão, convite de ingresso e chave determinística de pendência. CUE verifica campos, enums fechados, nulabilidade, condicionais (justificativa de exceção, autoria do aceite, modo de ingresso, contexto global) e limites locais; ingresso ordinário por vaga, exceção nominal acima da capacidade, proibição de editar capacidade abaixo da ocupação, unicidade concorrente de pendência, contador, remoção/reingresso, arquivamento somente leitura, idempotência M7 e composição M9 são verificados em Alloy; a canonicalização determinística da pendência (HMAC) e a proveniência por Rust. CUE não prova concorrência temporal.
\begin{description}
\item[id\_professor] Tipo: string. Obrigatório: sim. Aceita nulo: não.
SQL: TEXT.
Professor dono da turma; verificado no servidor.
\item[id\_materia] Tipo: string. Obrigatório: sim. Aceita nulo: não.
SQL: TEXT.
Matéria existente referenciada; não é recriada.
\item[nome\_turma] Tipo: string. Obrigatório: sim. Aceita nulo: não.
SQL: TEXT.
Nome de apresentação; no máximo 100 caracteres.
\item[ano] Tipo: int. Obrigatório: sim. Aceita nulo: não.
SQL: INTEGER.
Ano letivo inteiro.
\item[semestre] Tipo: int. Obrigatório: sim. Aceita nulo: não.
SQL: INTEGER.
Exclusivamente 1 ou 2.
\item[capacidade] Tipo: int. Obrigatório: sim. Aceita nulo: não.
SQL: INTEGER.
Máximo ordinário; só edição abaixo da ocupação é proibida (HQ-M11-001 = A).
\item[codigo\_turma] Tipo: string. Obrigatório: sim. Aceita nulo: não.
SQL: TEXT.
Código único gerado no servidor, até 20 caracteres; reservado em Chaves\_Unicas e não liberado ao arquivar.
\item[status\_turma] Tipo: enum. Obrigatório: sim. Aceita nulo: não.
SQL: ENUM.
Valores: Ativo, Arquivada.
\item[versao] Tipo: int. Obrigatório: sim. Aceita nulo: não.
SQL: INTEGER.
Detecta edição concorrente de metadados; inclusão/remoção de membro não incrementa.
\item[qtd\_alunos] Tipo: int. Obrigatório: sim. Aceita nulo: não.
SQL: INTEGER.
Contador transacional do vínculo ativo; pode exceder capacidade por exceção nominal válida.
\item[id\_aluno] Tipo: string. Obrigatório: sim. Aceita nulo: não.
SQL: TEXT.
UID do aluno no vínculo canônico.
\item[id\_turma] Tipo: string. Obrigatório: sim. Aceita nulo: sim.
SQL: TEXT.
Turma do vínculo/convite; nulo identifica convite global.
\item[email] Tipo: string. Obrigatório: sim. Aceita nulo: não.
SQL: TEXT.
E-mail normalizado do convite; no máximo 150 caracteres.
\item[token\_hash] Tipo: string. Obrigatório: sim. Aceita nulo: não.
SQL: TEXT.
SHA-256 do token aleatório de 32 bytes; nunca em claro.
\item[status\_convite] Tipo: enum. Obrigatório: sim. Aceita nulo: não.
SQL: ENUM.
Valores: pendente, aceitado, expirado.
\item[exceder\_capacidade] Tipo: bool. Obrigatório: sim. Aceita nulo: não.
SQL: BOOLEAN.
Exceção nominal do professor dono; justificativa obrigatória.
\item[modo\_ingresso] Tipo: enum. Obrigatório: sim. Aceita nulo: sim.
SQL: ENUM.
Valores: CODIGO, CONVITE.
\item[tipo\_evento] Tipo: enum. Obrigatório: sim. Aceita nulo: não.
SQL: ENUM.
Valores: inclusao\_aluno, exclusao\_aluno.
\item[aceitado\_por] Tipo: string. Obrigatório: sim. Aceita nulo: sim.
SQL: TEXT.
UID do aluno que aceitou; preservado no histórico.
\end{description}

% Gerado por lcqui-spec-doc; não editar.
\subsection*{Evidência formal M11 --- turma, matrícula e convite}
CUE verifica shapes de turma, vínculo canônico Aluno--Turma, espelho mínimo de consulta, evento de inclusão/exclusão, convite e chave de pendência, com enums fechados, nulabilidade, condicionais (justificativa de exceção, autoria do aceite, modo de ingresso, contexto global) e limites locais. Alloy verifica a unicidade e a reserva permanente de codigo\_turma, o arquivamento e desarquivamento preservando ID/membros/código/histórico e bloqueando escrita acadêmica enquanto arquivada, o ingresso ordinário apenas com vaga estrita em turma Ativa, a exceção nominal válida que pode exceder a capacidade, a proibição de editar capacidade abaixo da ocupação (HQ-M11-001 = A), a coerência vínculo canônico--contador--espelho, a remoção e o reingresso por convite, a unicidade de pendência por (e-mail, contexto) com distinção GLOBAL/turma, expiração, reenvio, terminalidade e preservação de histórico, a idempotência M7 (retry não duplica; reuso incompatível não herda) e a composição M9 (revogação/ownership impedem commit tardio; espelho não autoriza). Rust valida proveniência, ordem exata do receipt e as regras determinísticas de capacidade/exceção e canonicalização de e-mail. A evidência não certifica Firebase, Auth, envio de e-mail, HMAC concreto, Firestore/Storage Rules, backend, índices nem concorrência sob carga.
\begin{description}
\item[M11-INV-001] CodigoUnicoPorTurma: UNSAT.\newline Escopo: \texttt{check CodigoUnicoPorTurma for 6}.
\item[M11-INV-002] CriarTurmaReservaCodigo: UNSAT.\newline Escopo: \texttt{check CriarTurmaReservaCodigo for 6}.
\item[M11-INV-003] ArquivarPreservaMembros: UNSAT.\newline Escopo: \texttt{check ArquivarPreservaMembros for 6}.
\item[M11-INV-004] ArquivarNaoLiberaCodigo: UNSAT.\newline Escopo: \texttt{check ArquivarNaoLiberaCodigo for 6}.
\item[M11-INV-005] ArquivadaBloqueiaIngresso: UNSAT.\newline Escopo: \texttt{check ArquivadaBloqueiaIngresso for 6}.
\item[M11-INV-006] ArquivadaBloqueiaAceite: UNSAT.\newline Escopo: \texttt{check ArquivadaBloqueiaAceite for 6}.
\item[M11-INV-007] DesarquivarRestauraAtivo: UNSAT.\newline Escopo: \texttt{check DesarquivarRestauraAtivo for 6}.
\item[M11-INV-008] IngressoOrdinarioExigeVaga: UNSAT.\newline Escopo: \texttt{check IngressoOrdinarioExigeVaga for 6}.
\item[M11-INV-009] SemVagaNaoHaIngressoOrdinario: UNSAT.\newline Escopo: \texttt{check SemVagaNaoHaIngressoOrdinario for 6}.
\item[M11-INV-010] ExcecaoExigeConviteJustificado: UNSAT.\newline Escopo: \texttt{check ExcecaoExigeConviteJustificado for 6}.
\item[M11-INV-011] ExcecaoNaoEhBypassGenerico: UNSAT.\newline Escopo: \texttt{check ExcecaoNaoEhBypassGenerico for 6}.
\item[M11-INV-012] EdicaoCapacidadeNaoAbaixoOcupacao: UNSAT.\newline Escopo: \texttt{check EdicaoCapacidadeNaoAbaixoOcupacao for 6}.
\item[M11-INV-013] TransicoesPreservamContador: UNSAT.\newline Escopo: \texttt{check TransicoesPreservamContador for 6}.
\item[M11-INV-014] RemocaoPreservaHistorico: UNSAT.\newline Escopo: \texttt{check RemocaoPreservaHistorico for 6}.
\item[M11-INV-015] RemocaoBloqueiaCodigo: UNSAT.\newline Escopo: \texttt{check RemocaoBloqueiaCodigo for 6}.
\item[M11-INV-016] RemocaoAtualizaEspelho: UNSAT.\newline Escopo: \texttt{check RemocaoAtualizaEspelho for 6}.
\item[M11-INV-017] EspelhoNaoSobreviveRemocao: UNSAT.\newline Escopo: \texttt{check EspelhoNaoSobreviveRemocao for 6}.
\item[M11-INV-018] EspelhoEhProjecao: UNSAT.\newline Escopo: \texttt{check EspelhoEhProjecao for 6}.
\item[M11-INV-019] PendenciaUnicaPorEmailContexto: UNSAT.\newline Escopo: \texttt{check PendenciaUnicaPorEmailContexto for 6}.
\item[M11-INV-020] AceitarTurmaCriaVinculo: UNSAT.\newline Escopo: \texttt{check AceitarTurmaCriaVinculo for 6}.
\item[M11-INV-021] AceitarGlobalNaoMatricula: UNSAT.\newline Escopo: \texttt{check AceitarGlobalNaoMatricula for 6}.
\item[M11-INV-022] AceitarConsomeUmaVez: UNSAT.\newline Escopo: \texttt{check AceitarConsomeUmaVez for 6}.
\item[M11-INV-023] ExpiradoNaoOcupaPendencia: UNSAT.\newline Escopo: \texttt{check ExpiradoNaoOcupaPendencia for 6}.
\item[M11-INV-024] NovoConvitePreservaHistorico: UNSAT.\newline Escopo: \texttt{check NovoConvitePreservaHistorico for 6}.
\item[M11-INV-025] ReenvioMantemDocumento: UNSAT.\newline Escopo: \texttt{check ReenvioMantemDocumento for 6}.
\item[M11-INV-026] PendenciaGlobalDistinta: UNSAT.\newline Escopo: \texttt{check PendenciaGlobalDistinta for 6}.
\item[M11-INV-027] RevogacaoImpedeCommit: UNSAT.\newline Escopo: \texttt{check RevogacaoImpedeCommit for 6}.
\item[M11-INV-028] OwnershipErradoNaoGerencia: UNSAT.\newline Escopo: \texttt{check OwnershipErradoNaoGerencia for 6}.
\item[M11-INV-029] InativoNaoGerencia: UNSAT.\newline Escopo: \texttt{check InativoNaoGerencia for 6}.
\item[M11-INV-030] UsuarioInativoNaoLeTurma: UNSAT.\newline Escopo: \texttt{check UsuarioInativoNaoLeTurma for 6}.
\item[M11-INV-031] RetryNaoDuplicaFato: UNSAT.\newline Escopo: \texttt{check RetryNaoDuplicaFato for 6}.
\item[M11-INV-032] ReusoIncompativelNaoHerda: UNSAT.\newline Escopo: \texttt{check ReusoIncompativelNaoHerda for 6}.
\item[M11-INV-033] TransicoesPreservamCoerencia: UNSAT.\newline Escopo: \texttt{check TransicoesPreservamCoerencia for 6}.
\item[M11-WIT-034] WitnessTurmaAtiva: SAT.\newline Escopo: \texttt{run WitnessTurmaAtiva for 4}.
\item[M11-WIT-035] WitnessCriarTurma: SAT.\newline Escopo: \texttt{run WitnessCriarTurma for 4}.
\item[M11-WIT-036] WitnessArquivarComMembros: SAT.\newline Escopo: \texttt{run WitnessArquivarComMembros for 6}.
\item[M11-WIT-037] WitnessDesarquivar: SAT.\newline Escopo: \texttt{run WitnessDesarquivar for 6}.
\item[M11-WIT-038] WitnessIngressoOrdinario: SAT.\newline Escopo: \texttt{run WitnessIngressoOrdinario for 6}.
\item[M11-WIT-039] WitnessExcecaoAcimaCapacidade: SAT.\newline Escopo: \texttt{run WitnessExcecaoAcimaCapacidade for 6}.
\item[M11-WIT-040] WitnessCorridaUltimaVaga: SAT.\newline Escopo: \texttt{run WitnessCorridaUltimaVaga for 8}.
\item[M11-WIT-041] WitnessRemocao: SAT.\newline Escopo: \texttt{run WitnessRemocao for 6}.
\item[M11-WIT-042] WitnessReingressoPorConvite: SAT.\newline Escopo: \texttt{run WitnessReingressoPorConvite for 8}.
\item[M11-WIT-043] WitnessConviteGlobal: SAT.\newline Escopo: \texttt{run WitnessConviteGlobal for 4}.
\item[M11-WIT-044] WitnessAceitarConviteTurma: SAT.\newline Escopo: \texttt{run WitnessAceitarConviteTurma for 6}.
\item[M11-WIT-045] WitnessReenvioPendente: SAT.\newline Escopo: \texttt{run WitnessReenvioPendente for 4}.
\item[M11-WIT-046] WitnessNovoConviteAposTerminalidade: SAT.\newline Escopo: \texttt{run WitnessNovoConviteAposTerminalidade for 6}.
\item[M11-WIT-047] WitnessExpiracao: SAT.\newline Escopo: \texttt{run WitnessExpiracao for 4}.
\item[M11-WIT-048] WitnessEspelhoProjecao: SAT.\newline Escopo: \texttt{run WitnessEspelhoProjecao for 6}.
\item[M11-WIT-049] WitnessRevogacaoImpedeCommit: SAT.\newline Escopo: \texttt{run WitnessRevogacaoImpedeCommit for 6}.
\item[M11-WIT-050] WitnessEdicaoCapacidadeValida: SAT.\newline Escopo: \texttt{run WitnessEdicaoCapacidadeValida for 6}.
\end{description}
UNSAT para check indica ausência de contraexemplo no escopo declarado; SAT para run indica testemunha encontrada. A evidência não certifica a implementação Firebase/backend.

\par\endgroup

\subsection*{Limites}
Esta evidência prova somente o modelo formal no escopo declarado e \textbf{não}
certifica \texttt{functions/}\allowbreak\texttt{src/}\allowbreak\texttt{turmas.ts}, o frontend, \texttt{firestore.rules},
\texttt{storage.rules}, Firebase, Auth, Admin SDK, App Check, envio real de
e-mail, o HMAC concreto, transações reais sob carga, índices nem a
infraestrutura em produção. A prova de consistência do modelo não é homologação
da implementação real.
\fussy
